\Needspace{8\baselineskip}
\subsection{A signed three-parameter argument forcing sphericality}
\label{subsec:c3-signed-general}

\begin{proposition}\label{prop:c3-signed-general}
The rings $R7\_3\_783$ and $R7\_3\_1015$ have no categorification
over $\mathbb C$.
\end{proposition}
\begin{proof}
We show that every hypothetical categorification would be spherical;
Proposition~\ref{prop:pivotal-dimension-block} then gives a contradiction.
The two source rings have basis
$1,g,g^2,E,\rho_0,\rho_1,\rho_2$, where
$g^3=1$, $E^*=E$, $\rho_0^*=\rho_0$, and $\rho_i=g^i\rho_0$.
Put $G=1+g+g^2$ and $S=\rho_0+\rho_1+\rho_2$. Their rules are
\[
\begin{gathered}
 gE=Eg=E,\qquad \rho_0g=g^r\rho_0,\\
 E^2=G+3E+S,\qquad E\rho_i=\rho_iE=E+2S,
 \qquad \rho_0^2=1+2E+2S,
\end{gathered}
\]
where $r=-1$ for $R7\_3\_783$ and $r=1$ for $R7\_3\_1015$.
These relations generate all products. In census order one takes
$g=X_1$, $\rho_0=X_6$, and respectively $E=X_4,X_5$; the verifier
compares all $343$ structure constants.

Sphericalize a hypothetical category and pass to the signed quotient.
Choose the lift of $g$ with $g^3=1$, choose positive-second-indicator
lifts of $E$ and $\rho_0$, and define the other lifts by
$\rho_i=g^i\rho_0$. Since $gE$ has the same underlying simple as $E$,
its signed coefficient is $\pm1$. Its cube is one, so $gE=E$.
The same argument gives $Eg=E$ and
$\rho_0g=g^r\rho_0$.

The coefficient of $\rho_0$ in $E^2$ is one: its original multiplicity
is one and the positive-indicator square-channel congruence is modulo
four. Left multiplication by $g$ fixes $E^2$, forcing the other two
orbit coefficients also to be one. Frobenius reciprocity determines
the three invertible coefficients. Hence
\[
 E^2=G+bE+S,\qquad b\in\{-1,3\}.
\]
The element $E\rho_0$ is fixed by left multiplication by $g$, and
$E\rho_i=E\rho_0$. Its support and Frobenius reciprocity give
\[
 E\rho_i=E+aS,\qquad
 \rho_0^2=1+aE+\sum_{i=0}^2c_i\rho_i,
\]
where $a\in\{-2,2\}$ by the square-channel congruence, and
$|c_i|\leq2$. Put $c=\sum_i c_i\leq6$.
Multiplying the second identity by $G$ gives
$S\rho_0=G+3aE+cS$. Comparing the coefficients of $S$ in
$E^2\rho_0=E(E\rho_0)$ yields
\[
 ba+c=3a^2.
\]
The choices $(a,b)=(-2,-1),(-2,3),(2,-1)$ require respectively
$c=10,18,14$, all impossible. Thus $a=2$, $b=3$, and $c=6$,
forcing every $c_i=2$.

All displayed signed products are therefore the original positive
products. The $g$-relations generate the remaining products, so the
chosen seven lifts are closed under tensor product and duality. They
form a spherical fusion subcategory of the sphericalization. The
forgetful functor restricts to a tensor equivalence with the original
category, which is therefore spherical, as required.
\end{proof}

The signed argument is purely theoretical; the companion verification
checks only identification with the source tensors and elementary
arithmetic. In particular it does not presume that based-ring
symmetries furnish a pivotal structure. The remaining ring in
Proposition~\ref{prop:pivotal-dimension-block}, $R7\_3\_695$, is
excluded generally by Proposition~\ref{prop:r695-general} below.
